\chapter{Contextual Analysis}
\label{chap:six}

The empirical analysis of speech patterns, self vs other, and general topics provides great insight into how Bosnian politicians speak, but this alone cannot explain how their speech shapes ontological (in)security. This chapter contextualizes and expands upon the empirical findings from the text classification, sentiment analysis, and topic modelling in order to greater understand what Bosnian politicians view as the most prominent issues in \ac{BiH} and how they address them.

\section{Dysfunction in BiH}

A prominent theme across the 700 segments that were manually coded was the dysfunction of \ac{BiH}'s government especially regarding the three constituent peoples. This can be seen in quotes such as, 
\begin{quote}
    Here are the fundamental questions whether the will of one nation, another or third will be respected, and here is really the question---not the problem of the existence of three peoples, but in ... that we are not ready to admit it. And the third thing I want to do is we have a chance to get some lessons out of that 17-year-old post-Dayton fissure (\citealp{mochtak2022}; Translated via OPUS-MT).
\end{quote}
\ac{BiH} 

% Discuss narratives around Srebrenica & Dayton

